\documentclass[11pt,a4paper]{article}

\usepackage[margin=1in]{geometry}
\usepackage{hyperref}
\usepackage{enumitem}
\usepackage{booktabs}
\usepackage{tabularx}
\usepackage{xcolor}
\usepackage{mdframed}

\hypersetup{colorlinks=true, linkcolor=blue!70!black, urlcolor=blue!70!black}

\title{ERC Indoor Navigation Plan\\[0.3em]
\large System choice, module roles, and implementation plan}
\author{Project Working Draft}
\date{\today}

\begin{document}

\maketitle
\begin{center}
\fbox{\parbox{0.9\textwidth}{
\textbf{HISTORICAL DOCUMENT} \\[4pt]
Implementation plan from March 2026. Most of this plan was executed, but the
final result lives in the actual runtime code and later indoor narrative
documents. \\[4pt]
\textit{For the current system, see: \texttt{live\_indoor\_runtime\_story.tex},
\texttt{CLAUDE.md}}
}}
\end{center}
\vspace{1em}

\section{Read This First}

This document is written so that a teammate can understand the indoor system
without needing prior context from earlier discussions or code exploration.

\begin{mdframed}
\textbf{Short version:} we are solving ERC indoor navigation in a \textbf{known
corridor}. The system we will use is:

\begin{itemize}[nosep]
  \item a \textbf{CosPlace-based visual place recognition pipeline} to recognize locations,
  \item a \textbf{graph of reference images and recorded transitions} to represent the corridor,
  \item \textbf{SuperPoint/SuperGlue-style geometric verification} to reject weak matches,
  \item \textbf{graph planning} to choose the next nearby target place,
  \item a \textbf{validated online local controller} to drive from the current view to that nearby target image,
  \item a separate \textbf{safety and checkpoint verification layer}.
\end{itemize}
\end{mdframed}


\section{What Problem Are We Solving?}

The indoor ERC task is \textbf{image-goal navigation}.

That means:

\begin{itemize}[nosep]
  \item the robot is \textbf{not} given GPS coordinates,
  \item the robot is told to go to places that are represented by \textbf{images},
  \item the robot may have to visit \textbf{multiple checkpoints in sequence},
  \item the robot runs through the EarthRover / FrodoBots SDK over 4G,
  \item the robot sees only a few frames per second and commands arrive with delay.
\end{itemize}

So the real challenge is:

\begin{center}
\fbox{\parbox{0.9\textwidth}{
From a delayed low-FPS camera stream, figure out \textbf{where the robot is},
decide \textbf{which place it should go to next}, move there safely, and confirm
that the required checkpoint was actually reached.
}}
\end{center}


\section{Why Our Situation Is Better Than a General Indoor Navigation Problem}

This is important.

We are \textbf{not} dealing with a completely new and unseen environment.

We now know:

\begin{itemize}[nosep]
  \item we will test in the same corridor repeatedly,
  \item the competition corridor/environment family will also be the same one,
  \item we can collect data there in advance,
  \item we can tune our system for that exact corridor.
\end{itemize}

This changes everything.

\subsection*{What this allows us to do}

\begin{itemize}[nosep]
  \item build a database of images from that exact corridor,
  \item choose good reference locations manually,
  \item tune descriptor thresholds for that corridor,
  \item tune image cropping for white walls / floor reflections / bright doors,
  \item map each checkpoint image to known graph nodes ahead of time,
  \item accept some environment-specific overfitting on purpose.
\end{itemize}

\subsection*{What still does \underline{not} become easy}

\begin{itemize}[nosep]
  \item low FPS,
  \item 500 ms-ish latency,
  \item pedestrians,
  \item repeated corridor appearance,
  \item bright doorway / reflective floor effects,
  \item stable control on the real robot.
\end{itemize}


\section{What Existing Work We Already Have}

We already have code for a \textbf{CosPlace + SuperGlue corridor localization and graph-navigation pipeline}.

This work was focused on the \textbf{perception and planning} side of the problem.
It was not yet a full EarthRover runtime system, but it already implemented the
core logic for recognizing places and building a navigation graph from corridor data.

\subsection*{What that pipeline was designed to do}

The pipeline does the following:

\begin{enumerate}[nosep]
  \item record images while moving through the corridor,
  \item represent each place using a visual descriptor,
  \item connect those places into a graph,
  \item when a new query image comes in, find the most similar places in the graph,
  \item verify the match using geometry,
  \item then use the graph to reason about where the robot is and how to move.
\end{enumerate}

\subsection*{What technical words mean in plain English}

\textbf{Visual descriptor}

A compact vector that represents what a place looks like.
If two images are visually similar, their descriptors should be similar.

\textbf{Retrieval}

Given the current camera frame, search the stored database and find the most
similar reference images.

\textbf{Graph}

A set of nodes and edges.
Here:

\begin{itemize}[nosep]
  \item a \textbf{node} means a saved reference image / place,
  \item an \textbf{edge} means two places are connected or reachable.
\end{itemize}

\textbf{Geometric verification}

After retrieval says ``these two images look similar,'' geometric verification
checks whether they are actually consistent in terms of visual keypoint matches,
instead of being a false match caused by a repetitive corridor.

\textbf{Action graph}

A directed graph that says not just which places are similar, but how recorded
motion connected one place to the next.

\textbf{MBRA}

A learned visual controller that takes recent observation images and a goal image
and predicts short-horizon motion.

In the MBRA paper/codebase, its primary role is the \textbf{short-horizon expert
used for relabeling or annotation}. The deployed-policy side of that project is
\textbf{LogoNav}.


\section{What The CosPlace + SuperGlue Code Already Does}

The file \texttt{baseline.py} already gives us a real
\textbf{perception/planning baseline}.

It now serves as a shared baseline script with a database-building workflow and
a query/localization workflow.

\subsection*{It already does these things}

\begin{itemize}[nosep]
  \item computes global visual descriptors for corridor images,
  \item builds a database of those descriptors,
  \item retrieves nearest visual matches for a query image,
  \item builds a place graph over image nodes,
  \item builds an action / connectivity graph from recorded route data,
  \item runs SuperPoint/SuperGlue-style geometric verification,
  \item visualizes graphs and candidate matches for debugging.
\end{itemize}

\subsection*{This means the current corridor pipeline already solves}

\begin{itemize}[nosep]
  \item how to represent places,
  \item how to compare a live image to stored places,
  \item how to connect those places into a graph,
  \item how to reject weak or fake matches,
  \item how to inspect what the system thinks is happening.
\end{itemize}

\subsection*{But it does \underline{not} yet solve}

\begin{itemize}[nosep]
  \item live EarthRover SDK runtime,
  \item real-time command loop,
  \item multi-checkpoint mission execution,
  \item selection and validation of the final online local controller,
  \item safety override,
  \item recovery behavior if localization becomes uncertain.
\end{itemize}

\begin{mdframed}
\textbf{Correct conclusion:} the \textbf{CosPlace + SuperGlue corridor pipeline}
should be treated as the actual baseline for \textbf{perception and planning},
but not as the full competition system.
\end{mdframed}


\section{What We Are Building Now}

We are building a full system around that \textbf{CosPlace + SuperGlue + graph}
backbone.

The complete system has five parts:

\begin{enumerate}[nosep]
  \item perception / localization,
  \item planning,
  \item online local control,
  \item safety,
  \item checkpoint verification and recovery.
\end{enumerate}


\section{Big Picture: Two Phases}

\subsection{Phase A: before competition day}

Before runtime, we prepare the environment model.

\begin{enumerate}[nosep]
  \item drive the robot or manually collect images through the corridor,
  \item save those images as reference images,
  \item compute CosPlace descriptors for them,
  \item build a graph of the corridor,
  \item attach recorded route/action information,
  \item decide which graph nodes correspond to likely checkpoint regions,
  \item test retrieval and SuperGlue-based verification on repeated runs.
\end{enumerate}

\subsection{Phase B: during runtime}

During competition/runtime, the robot repeatedly does this:

\begin{enumerate}[nosep]
  \item read current front-camera frame,
  \item localize in the corridor graph,
  \item plan path to the current checkpoint target node,
  \item choose the next nearby subgoal node,
  \item give that subgoal image to the chosen online controller,
  \item get a local control command,
  \item run safety checks,
  \item send safe command to the rover,
  \item check whether the current checkpoint has been reached,
  \item if yes, switch to the next checkpoint.
\end{enumerate}


\section{Each Stack Explained Clearly}

\subsection{1. Perception / Localization Stack}

\textbf{Job:} estimate where the robot is in the corridor graph.

\textbf{Input:}

\begin{itemize}[nosep]
  \item current live camera frame,
  \item reference database of corridor images and descriptors,
  \item previous localization state.
\end{itemize}

\textbf{What it does step by step:}

\begin{enumerate}[nosep]
  \item compute a CosPlace descriptor for the current frame,
  \item retrieve top matching reference images,
  \item run SuperGlue-based geometric verification on the best candidates,
  \item combine this with temporal continuity,
  \item output current best node estimate and confidence.
\end{enumerate}

\textbf{Output:}

\begin{itemize}[nosep]
  \item current node estimate,
  \item candidate nodes,
  \item confidence score,
  \item localization diagnostics.
\end{itemize}

\textbf{What team members should understand:}

This stack mainly reuses the existing CosPlace + SuperGlue code, but it must be
turned from notebook-style analysis code into a \textbf{live runtime module}.


\subsection{2. Planning Stack}

\textbf{Job:} choose the next place image the robot should move toward.

\textbf{Input:}

\begin{itemize}[nosep]
  \item current node estimate,
  \item graph of corridor nodes,
  \item current target checkpoint,
  \item mission progress state.
\end{itemize}

\textbf{What it does step by step:}

\begin{enumerate}[nosep]
  \item map the current checkpoint to a target node or target region,
  \item compute a path on the corridor graph from the current node to that target,
  \item choose a nearby node on that path as the current subgoal,
  \item keep that subgoal stable until there is enough evidence to switch,
  \item once the checkpoint is confirmed, move to the next checkpoint.
\end{enumerate}

\textbf{Output:}

\begin{itemize}[nosep]
  \item current graph path,
  \item next subgoal node,
  \item next subgoal image,
  \item checkpoint progression state.
\end{itemize}

\textbf{Important rule:}

The clean starting point is for planning to give the online controller a \textbf{nearby graph-node image}.
If someone proposes a stronger alternative, it should be evaluated against this baseline on the real corridor.

\subsection{3. Online Local Control Stack}

\textbf{Job:} move from the current camera view to the next nearby graph-node image.

\textbf{Input:}

\begin{itemize}[nosep]
  \item recent camera frames,
  \item current subgoal image from the planner,
  \item past command history,
  \item delay-related inputs if needed by the chosen controller.
\end{itemize}

\textbf{What it does step by step:}

\begin{enumerate}[nosep]
  \item collect the recent visual context,
  \item package the current subgoal image,
  \item run the chosen online controller,
  \item convert controller output into rover linear/angular commands,
  \item smooth or limit commands if needed.
\end{enumerate}

\textbf{Output:}

\begin{itemize}[nosep]
  \item proposed motion command.
\end{itemize}

\textbf{Important boundary:}

The online local controller is not the full navigation system.
It is only the short-horizon movement module between nearby nodes.

\textbf{Important clarification about MBRA / LogoNav}

\begin{itemize}[nosep]
  \item In the paper, \textbf{MBRA} is mainly the short-horizon expert used for relabeling.
  \item \textbf{LogoNav} is the deployed-policy side of that project.
  \item In this repo, the provided deployment scripts are for \textbf{LogoNav}, not for an indoor image-goal MBRA controller.
  \item LogoNav here is GPS/pose-conditioned, so it is not a drop-in indoor image-goal controller.
\end{itemize}

So one of the current engineering tasks is to determine what the actual online
local controller will be for this indoor stack.


\subsection{4. Safety Stack}

\textbf{Job:} answer the question ``Is the proposed motion safe enough to execute?''

\textbf{Input:}

\begin{itemize}[nosep]
  \item current camera frame,
  \item proposed local-control command,
  \item optional depth/clearance estimate,
  \item optional pedestrian detector,
  \item localization confidence.
\end{itemize}

\textbf{What it does step by step:}

\begin{enumerate}[nosep]
  \item check whether the command looks safe,
  \item slow or stop if a person is ahead,
  \item slow down if localization confidence is weak,
  \item veto dangerous commands,
  \item send only a safe final command.
\end{enumerate}

\textbf{Output:}

\begin{itemize}[nosep]
  \item final command sent to the robot,
  \item safety state such as OK / slow / stop.
\end{itemize}


\subsection{5. Checkpoint Verification and Recovery Stack}

\textbf{Checkpoint verification job:} decide whether the required checkpoint has actually been reached.

\textbf{What it should do:}

\begin{itemize}[nosep]
  \item require repeated agreement over multiple frames,
  \item use geometric verification for final confirmation,
  \item avoid declaring success from one accidental similar corridor view.
\end{itemize}

\textbf{Recovery job:} answer the question ``What do we do if we are not sure where we are?''

\textbf{What it should do:}

\begin{itemize}[nosep]
  \item stop or slow down instead of continuing blindly,
  \item try relocalizing from recent frames,
  \item optionally perform a small search motion,
  \item resume only after confidence improves.
\end{itemize}


\section{How All Stacks Connect Together}

\begin{center}
\fbox{\parbox{0.94\textwidth}{
Camera frame $\rightarrow$
Perception says current node/confidence $\rightarrow$
Planner chooses path and nearby subgoal $\rightarrow$
online controller proposes local command $\rightarrow$
Safety accepts/modifies/rejects command $\rightarrow$
Robot moves $\rightarrow$
Checkpoint verifier decides whether to advance mission state
}}
\end{center}

\subsection*{In one sentence}

\begin{itemize}[nosep]
  \item Perception says where we are in the corridor graph.
  \item Planning says which nearby graph node we should move toward next.
  \item The online controller says how to move locally toward that node image.
  \item Safety says whether that movement is safe enough to execute.
  \item Checkpoint verification says whether the current checkpoint is complete.
\end{itemize}


\section{Why We Are Not Using Other Approaches As the Main Backbone}

\subsection*{Why not pure end-to-end policy?}

Because we already have a more debuggable and corridor-specific system based on
CosPlace retrieval, geometric verification, and graph planning. A pure end-to-end
policy would hide localization, routing, and failure recovery inside one black box.

\subsection*{Why not ORB-SLAM3 as the main backbone?}

Because even in a known corridor, our real runtime still has:

\begin{itemize}[nosep]
  \item low FPS,
  \item latency,
  \item visually repetitive corridor structure,
  \item reflective floor,
  \item monocular RGB limitations.
\end{itemize}

ORB-SLAM3 can still be tested, but it should not replace a corridor-specific
CosPlace + graph baseline unless it proves clearly better on our real data.

\subsection*{Why not retrieval-only without a local controller?}

Because retrieval tells us \textbf{which place looks similar}, but it does not
by itself provide a strong short-horizon controller under delay and low frame rate.


\section{What Becomes Easier Because the Corridor Is Known}

\begin{itemize}[nosep]
  \item node spacing can be tuned on the exact route,
  \item descriptor preprocessing can be corridor-specific,
  \item bright-door and reflective-floor cases can be tested directly,
  \item checkpoint thresholds can be calibrated on real repeated data,
  \item ambiguous corridor segments can be identified in advance,
  \item some overfitting is acceptable because this is the actual target environment.
\end{itemize}

\subsection*{Examples of acceptable overfitting}

\begin{itemize}[nosep]
  \item cropping ceiling/floor bands,
  \item putting more nodes near doors or signs,
  \item checkpoint-specific verification thresholds,
  \item graph pruning for impossible transitions.
\end{itemize}


\section{What Can Still Fail}

Even with a known corridor, these are the main remaining risks:

\begin{enumerate}[nosep]
  \item The chosen online controller fails to follow a nearby subgoal cleanly.
  \item Node switching becomes unstable in visually repetitive areas.
  \item Checkpoint verification declares success too early.
  \item Bright doorway / reflective floor causes bad visual matching.
  \item Pedestrians block the visual cues we need.
  \item Low FPS and latency cause overshoot or slow correction.
\end{enumerate}


\section{Who Should Work On What}

\begin{tabularx}{\textwidth}{p{0.17\textwidth} p{0.31\textwidth} X}
\toprule
\textbf{Area} & \textbf{Main responsibility} & \textbf{What they should deliver} \\
\midrule
Perception & Turn the CosPlace + SuperGlue corridor code into a live localization module & Current-node estimate, confidence score, candidate matches, runtime localization API \\
Planning & Build path query and sequential checkpoint logic on the graph & Target-node lookup, path planner, subgoal selector, mission state logic \\
Control & Select and validate the online local controller on the rover & Controller input pipeline, local command generation, smoothing / limits \\
Safety & Prevent unsafe motion at runtime & Slow/stop logic, pedestrian-aware checks, confidence-aware speed limiting \\
Checkpoint / Recovery & Decide when a checkpoint is really reached and what to do when uncertain & Goal verifier, recovery state logic, relocalization behavior \\
Integration & Connect everything end-to-end with the SDK & One full runnable ERC indoor loop \\
Testing & Repeated corridor experiments and tuning & Metrics, debug plots, threshold tuning, failure-case reports \\
\bottomrule
\end{tabularx}


\section{What We Should Do Next In Order}

\begin{enumerate}[nosep]
  \item Reuse the existing CosPlace + SuperGlue corridor code directly.
  \item Clean it up into reusable modules instead of notebook-style code.
  \item Validate live localization on the exact corridor.
  \item Determine which online local controller is actually deployable.
  \item Validate that controller on nearby graph-node subgoals.
  \item Integrate planner $\rightarrow$ subgoal $\rightarrow$ controller.
  \item Add safety and checkpoint verification.
  \item Add recovery behavior.
  \item Then run full multi-checkpoint tests.
\end{enumerate}

\textbf{Lower-priority directions right now:}

\begin{itemize}[nosep]
  \item rebuilding the CosPlace + graph backbone from scratch,
  \item replacing the full stack before the current baseline is validated end-to-end,
  \item making SLAM the default backbone before it proves better on our corridor,
  \item training a big new policy before the known-corridor baseline works.
\end{itemize}

Improvements are welcome, but the standard should be clear:
\textbf{a change is worth adopting if it makes the real robot more reliable on the real ERC indoor route.}


\section{Final Team Verdict}

\begin{mdframed}
\begin{itemize}[nosep]
  \item The \textbf{CosPlace + SuperGlue corridor-localization pipeline} is the real baseline for perception and planning.
  \item The main new work is \textbf{runtime integration around that baseline}.
  \item The \textbf{online local controller} still has to be selected and validated carefully.
  \item Safety, checkpoint verification, and recovery must be separate modules.
  \item The biggest risk is \textbf{runtime stability on the real robot}, not lack of architecture ideas.
\end{itemize}
\end{mdframed}

\end{document}
